\section*{Bab \ref{ch:g2:water-on-earth} --- Air di Bumi}

\begin{solution}{exo:g2:water-on-earth:1}
\omterm{def:g2:water-on-earth:salt-fresh}{Air asin}: \omterm{def:g2:water-on-earth:ocean}{laut}, \omterm{def:g2:water-on-earth:ocean}{samudra}. \omterm{def:g2:water-on-earth:salt-fresh}{Air tawar}: \omterm{def:g1:landscapes-around-us:water}{sungai}, hujan, \omterm{def:g1:landscapes-around-us:water}{danau} di pegunungan, es
di \omterm{def:g1:landscapes-around-us:landform}{gunung}.
\end{solution}

\begin{solution}{exo:g2:water-on-earth:2}
Karena sebagian besar permukaannya ditutupi air biru \omterm{def:g2:water-on-earth:ocean}{samudra}: dilihat dari
antariksa, Bumi tampak biru.
\end{solution}

\begin{solution}{exo:g2:water-on-earth:3}
$100 - 71 = 29$ persegi berupa daratan.
\end{solution}

\begin{solution}{exo:g2:water-on-earth:4}
\omterm{def:g2:water-on-earth:ocean}{Samudra} lebih besar: \omterm{def:g2:water-on-earth:ocean}{laut} adalah hamparan \omterm{def:g2:water-on-earth:salt-fresh}{air asin} yang lebih kecil, sering
di tepi \omterm{def:g2:water-on-earth:ocean}{samudra}.
\end{solution}

\begin{solution}{exo:g2:water-on-earth:5}
Tidak: air \omterm{def:g2:water-on-earth:ocean}{samudra} adalah \omterm{def:g2:water-on-earth:salt-fresh}{air asin}, dan meminumnya membuat tubuh kehilangan
air. Pelaut itu harus membawa \omterm{def:g2:water-on-earth:salt-fresh}{air tawar}.
\end{solution}

\begin{solution}{exo:g2:water-on-earth:6}
Misalnya: di \omterm{def:g1:landscapes-around-us:water}{sungai}, di \omterm{def:g1:landscapes-around-us:water}{danau}, di es dan salju di pegunungan dan di
Antarktika, di bawah tanah, di dalam air hujan.
\end{solution}

\begin{solution}{exo:g2:water-on-earth:7}
\omterm{def:g2:water-on-earth:salt-fresh}{Air asin}: 96 ember. \omterm{def:g2:water-on-earth:salt-fresh}{Air tawar}: $2 + 2 = 4$ ember.
(Sebenarnya sebagian air di bawah tanah juga asin; gambar itu menghitungnya
bersama \omterm{def:g2:water-on-earth:salt-fresh}{air tawar}.)
\end{solution}

\begin{solution}{exo:g2:water-on-earth:8}
\omterm{def:g2:water-on-earth:ocean}{Samudra} (96 ember), lalu es dan salju (2 ember), lalu \omterm{def:g1:landscapes-around-us:water}{sungai} (jauh kurang
dari satu ember).
\end{solution}

\begin{solution}{exo:g2:water-on-earth:9}
Air itu bisa meresap ke dalam tanah; bisa mengalir di atas tanah ke \omterm{def:g1:landscapes-around-us:water}{sungai}
kecil atau \omterm{def:g1:landscapes-around-us:water}{sungai}; bisa menguap kembali ke udara ketika matahari muncul.
\end{solution}

\begin{solution}{exo:g2:water-on-earth:10}
Garam tertinggal di dasar panci. Airnya sudah menguap: air itu pergi ke
udara sebagai uap air.
\end{solution}

\begin{solution}{exo:g2:water-on-earth:11}
Misalnya: tetes itu mengalir turun dari atap ke talang dan selokan, lalu ke
\omterm{def:g1:landscapes-around-us:water}{sungai} kecil dan \omterm{def:g1:landscapes-around-us:water}{sungai}; \omterm{def:g1:landscapes-around-us:water}{sungai} membawanya ke \omterm{def:g2:water-on-earth:ocean}{laut}; pada suatu hari yang
cerah tetes itu menguap dari \omterm{def:g2:water-on-earth:ocean}{laut}, naik ke udara, dan menjadi bagian dari
sebuah awan.
\end{solution}
